\section{Reduced homomorphisms and nilpotency ascent}
\label{section: relative reduced homomorphisms}

In this section, we introduce reduced and operator-reduced surjective homomorphisms and establish the nilpotency ascent theorem.


\subsection{Reduced and operator-reduced homomorphisms}
\label{subsection: definition of relative reduced homomorphisms}

Following \cite[Definition 3.1]{inoue_2013_quasitriviality_of_quandles_for_linkhomotopy} and \cite[Definition 7.1]{darne_2026_nilpotent_quandles} (cf.~\cite[Definition 4.1]{hughes_2011_link_homotopy_invariant_quandles}), a quandle $P$ is called \emph{reduced} if $s_{\varphi(x)}(x)=x$ for every $x\in P$ and every $\varphi\in\Inn(P)$.

We introduce a relative version of reduced quandles, as well as a slightly more general notion called \emph{operator-reduced} (cf.~\cite[Section 5]{hughes_2011_link_homotopy_invariant_quandles}).

\begin{dfn}
\label[dfn]{definition: reduced hom and operator reduced hom}
A surjective quandle homomorphism $f\colon P\twoheadrightarrow Q$ is called
\begin{enumerate}
    \item \emph{reduced} if $s_{\varphi(x)}(x)=x$, and
    \item \emph{operator-reduced} if $s_xs_{\varphi(x)}=s_{\varphi(x)}s_x$
\end{enumerate}
for every $x\in P$ and every $\varphi\in\relInn(f)$.
\end{dfn}

For the terminal homomorphism $P\twoheadrightarrow\{*\}$, the first condition is precisely the notion of a reduced quandle. Since $\relInn(f)\subseteq\Inn(P)$, every surjective homomorphism with reduced domain is reduced.

The two conditions are related to the first relative commutators.

\begin{prop}
\label[prop]{proposition: commutator interpretation of reducedness}
Let $f\colon P\twoheadrightarrow Q$ be a surjective quandle homomorphism.
For $x\in P$ and $\varphi\in\relInn(f)$, it holds that $s_{\varphi(x)}=[\varphi,s_x]s_x$ and $[s_x,s_{\varphi(x)}]=[[\varphi,s_x],s_x]^{-1}$.
Consequently, the following hold.
\begin{enumerate}
    \item\label{item: 1 nil} $f$ is $1$-nilpotent if and only if $[\varphi,s_x]=1$ for every $x$ and $\varphi$.
    \item\label{item: reduced} $f$ is reduced if and only if $[\varphi,s_x](x)=x$ for every $x$ and $\varphi$.
    \item\label{item: op reduced} $f$ is operator-reduced if and only if $[[\varphi,s_x],s_x]=1$ for every $x$ and $\varphi$.
\end{enumerate}
In particular, we have the implications $1$-nilpotent $\Rightarrow$ reduced $\Rightarrow$ operator-reduced and $2$-nilpotent $\Rightarrow$ operator-reduced.
\end{prop}

\begin{proof}
Since $\varphi\in\Inn(P)$, we have $s_{\varphi(x)}=\varphi s_x\varphi^{-1}$. Hence, we have $s_{\varphi(x)}=[\varphi,s_x]s_x$ and $[s_x,s_{\varphi(x)}]=[[\varphi,s_x],s_x]^{-1}$.
For \eqref{item: 1 nil}, since the symmetry maps generate $\Inn(P)$, the condition is equivalent to $[\Inn(P),\relInn(f)]=1$. For \eqref{item: reduced}, the first identity and $s_x(x)=x$ give $s_{\varphi(x)}(x)=[\varphi,s_x](x)$.
The equivalence in \eqref{item: op reduced} follows from the second identity.

For the last implication, since $[[\varphi,s_x],s_x]\in\relInnLCS{2}{f}$, \eqref{item: op reduced} shows that $2$-nilpotent homomorphisms are operator-reduced.
\end{proof}


\begin{prop}
\label[prop]{proposition: reduced Alexander homomorphisms}
In the Alexander setting of
\cref{example: relative Alexander lower central series}, the following
hold.
\begin{enumerate}
    \item\label{item: reduced Alexander homomorphisms}
    The homomorphism $f$ is reduced if and only if
    $(1-\sigma)D=0$, or equivalently, if and only if $f$ is
    $1$-nilpotent.
    \item\label{item: operator-reduced Alexander homomorphisms}
    The homomorphism $f$ is operator-reduced if and only if
    $(1-\sigma)^2D=0$, or equivalently, if and only if $f$ is
    $2$-nilpotent.
\end{enumerate}
\end{prop}

\begin{proof}
As observed in the proof of
\cref{example: relative Alexander lower central series},
we have $\varphi(x)-x\in D$ for every $x\in A$ and
$\varphi\in\relInn(f)$. On the other hand,
$L(D)\subseteq\relInn(f)$ by \cref{example: Alexander quandle}.
Hence,
\(
    \relInn(f)\cdot x=x+D
\)
for every $x\in A$.
Therefore, $f$ is reduced if and only if
$s_{x+d}(x)=x$ for every $x\in A$ and $d\in D$.
Since
\(
    s_{x+d}(x)=x+(1-\sigma)d,
\)
this is equivalent to $(1-\sigma)D=0$.
By \cref{example: relative Alexander lower central series},
this is further equivalent to $\Gamma_1(f)=1$, that is, to
$f$ being $1$-nilpotent.

Similarly, $f$ is operator-reduced if and only if
$s_x$ and $s_{x+d}$ commute for every $x\in A$ and $d\in D$.
Since
\(
    s_{x+d}=L_{(1-\sigma)d}s_x
\)
and conjugation by $s_x$ sends $L_u$ to $L_{\sigma(u)}$,
these two symmetry maps commute if and only if
\(
    \sigma((1-\sigma)d)=(1-\sigma)d,
\)
or equivalently, $(1-\sigma)^2d=0$.
Thus, $f$ is operator-reduced if and only if
$(1-\sigma)^2D=0$. Again by
\cref{example: relative Alexander lower central series},
this is equivalent to $\Gamma_2(f)=1$, that is, to
$f$ being $2$-nilpotent.
\end{proof}

\begin{prop}
\label[prop]{proposition: characterizations of reduced homomorphisms}
For a surjective quandle homomorphism $f\colon P\twoheadrightarrow Q$, the following conditions are equivalent.
\begin{enumerate}[label={\textup{(\roman*)}}]
    \item\label{item: f is reduced} $f$ is reduced.
    \item\label{item: orbit is trivial} Every $\relInn(f)$-orbit in $P$ is a trivial subquandle.
\end{enumerate}
In particular, if every fiber of $f$ is trivial, seen as a subquandle of $P$, then $f$ is reduced. The converse holds when $f$ is connected.
\end{prop}

\begin{proof}
\ref{item: orbit is trivial} $\Rightarrow$ \ref{item: f is reduced}: Consider an orbit $P'\coloneqq \relInn(f)\cdot x \subseteq P$. Since it is a trivial subquandle, $s_{\varphi(x)}\rvert_{P'}=\id_{P'}$ for all $\varphi\in \relInn(f)$. In particular, $s_{\varphi(x)}(x)=\id(x)=x$, and hence $f$ is reduced.
\ref{item: f is reduced}$\Rightarrow$ \ref{item: orbit is trivial}: if we take $y=\varphi(x)$ and $z=\psi(x)$ with $\varphi,\psi\in\relInn(f)$ from the orbit $\relInn(f)\cdot x$, then $s_y(z)=\psi(s_{\psi^{-1}\varphi(x)}(x))=\psi(x)=z$, which shows that $\relInn(f)\cdot x$ is a trivial subquandle.

The last assertions follow because every $\relInn(f)$-orbit is contained in a fiber, and the two coincide when $f$ is connected.
\end{proof}

In particular, covering homomorphisms are reduced, since we can easily verify that every fiber of a covering homomorphism is a trivial subquandle.

\begin{prop}
\label[prop]{proposition: first properties of reduced homomorphisms}
Let $f\colon P\to Q$, $f'\colon P'\to Q'$, and $g\colon Q\to R$ be surjective quandle homomorphisms. Put $h\coloneqq g\circ f$.
\begin{enumerate}
    \item\label{item: (operator-)reduced descent} If $h$ is (operator-)reduced, then both $f$ and $g$ are (operator-)reduced.
    \item\label{item: product of (operator-)reduced homs} If $f$ and $f'$ are (operator-)reduced, then $f\times f'$ is (operator-)reduced.
    \item\label{item: (operator-)reduced pullback} If $f$ is (operator-)reduced, then every pullback of $f$ along a surjective quandle homomorphism is (operator-)reduced.
    \item\label{item: connected (operator-)reduced pullback} If $f$ is connected and (operator-)reduced, then every pullback of $f$ along an arbitrary quandle homomorphism is (operator-)reduced.
\end{enumerate}
\end{prop}

\begin{proof}
\eqref{item: (operator-)reduced descent}:
Since $\relInn(f)\subseteq\relInn(h)$, if $h$ is reduced or operator-reduced, then so is $f$. Let $y\in Q$ and $\psi\in\relInn(g)$. Choose $x\in P$ with $f(x)=y$. By \cref{lemma: relative symmetry groups under a surjective factorization}, there is $\varphi\in\relInn(h)$ such that $f_*(\varphi)=\psi$.
If $h$ is reduced, then $s_{\psi(y)}(y) = f\bigl(s_{\varphi(x)}(x)\bigr) = f(x) = y$, so $g$ is reduced.
If $h$ is operator-reduced, then $[s_y,s_{\psi(y)}] = f_*([s_x,s_{\varphi(x)}])=1$. Hence, $g$ is operator-reduced.

\eqref{item: product of (operator-)reduced homs}:
Let $z=(x,x')\in P\times P'$ and $\varphi\in\relInn(f\times f')$. Through the injection $\Inn(P\times P') \hookrightarrow \Inn(P) \times \Inn(P')$ from \cref{proposition: inner automorphism groups of products}, we identify $\varphi=(\varphi_1,\varphi_2)$, where $\varphi_1\in\relInn(f)$ and $\varphi_2\in\relInn(f')$ are its projections. Hence we have $s_{\varphi(z)}=(s_{\varphi_1(x)},s_{\varphi_2(x')})$.
If $f$ and $f'$ are reduced, then we have $s_{\varphi(z)}(z)= \bigl(s_{\varphi_1(x)}(x), s_{\varphi_2(x')}(x')\bigr)=(x,x')=z$, so $f\times f'$ is reduced.
If $f$ and $f'$ are operator-reduced, then we have $s_zs_{\varphi(z)}=(s_x,s_{x'})(s_{\varphi_1(x)},s_{\varphi_2(x')}) = (s_{\varphi_1(x)},s_{\varphi_2(x')})(s_x,s_{x'}) = s_{\varphi(z)}s_z$, so $f\times f'$ is operator-reduced.

\eqref{item: (operator-)reduced pullback}:
Consider a pullback square of $f$ along a surjective homomorphism:
\[\begin{tikzcd}
P' \arrow[r, two heads, "u"] \arrow[d, two heads, "f'"']
    & P \arrow[d, two heads, "f"]\\
Q' \arrow[r, two heads, "v"']
    & Q\rlap{.} \arrow[phantom]{ul}[very near end]{\lrcorner}
\end{tikzcd}\]
Let $z\in P'$ and $\psi\in\relInn(f')$. Then $u_*(\psi)\in\relInn(f)$. Moreover $f'_*(\psi)=1$ gives $f's_{\psi(z)}=s_{f'\psi(z)}f'=s_{f'(z)} f'$.
If $f$ is reduced, then we have
\[ u\bigl(s_{\psi(z)}(z)\bigr) = s_{u\psi(z)}(u(z)) = s_{u_*(\psi)(u(z))}(u(z)) = u(z), \]
as well as $f'(s_{\psi(z)}(z)) = s_{f'(z)}(f'(z)) = f'(z)$. Since $(u,f')$ is injective by the pullback property, $s_{\psi(z)}(z)=z$, so $f'$ is reduced.
If $f$ is operator-reduced, then we have
\[ us_zs_{\psi(z)} = s_{u(z)}s_{u\psi(z)} u = s_{u(z)}s_{u_*(\psi)(u(z))}u = s_{u_*(\psi)(u(z))}s_{u(z)}u = u s_{\psi(z)} s_z, \]
as well as $f's_zs_{\psi(z)} = s_{f'(z)}s_{f'(z)} f' = f's_{\psi(z)} s_z$. Hence the universality of pullback yields $s_zs_{\psi(z)}= s_{\psi(z)} s_z$, so $f'$ is operator-reduced.

\eqref{item: connected (operator-)reduced pullback}:
Suppose first that $f$ is reduced. By
\cref{proposition: characterizations of reduced homomorphisms},
every fiber of $f$ is trivial. Every fiber of a pullback of $f$
is isomorphic to a fiber of $f$, and hence is trivial. The assertion
follows from the same proposition.

Suppose that $f$ is operator-reduced.
If $x,y\in P$ satisfy $f(x)=f(y)$, then the connectedness of $f$
gives some $\varphi\in\relInn(f)$ such that $y=\varphi(x)$.
Hence $[s_x,s_y]=1$.
Consider an arbitrary pullback $f'\colon P'\to Q'$ of $f$.
For $z=(x,q')\in P'$ and $\psi\in\relInn(f')$, write
$\psi(z)=(y,q')$. Then $f(x)=f(y)$, and hence $[s_x,s_y]=1$.
Since $s_z=(s_x,s_{q'})|_{P'}$ and $s_{\psi(z)}=(s_y,s_{q'})|_{P'}$,
we obtain $[s_z,s_{\psi(z)}]=1$. Thus, $f'$ is operator-reduced.
\end{proof}






\subsection{Universal (operator-)reduced factor}
\label{subsection: universal relative reduced quotient}


Let $f\colon P\twoheadrightarrow Q$ be a surjective quandle homomorphism. 
Let $\redCong{f}$ be the quandle congruence generated by $s_{\varphi(x)}(x)\mathrel{\redCong{f}}x$ for $x\in P$ and $\varphi\in\relInn(f)$. 
Let $\opredDef{f}$ be the normal subgroup of $\Inn(P)$ generated by $[s_x,s_{\varphi(x)}]$ for $x\in P$ and $\varphi\in\relInn(f)$.
The first relation lies in the kernel congruence $\Eq(f)$, and \cref{proposition: commutator interpretation of reducedness} gives $\opredDef{f}\subseteq\relInnLCS{2}{f}\subseteq \relInn(f)$.

\begin{prop}
\label[prop]{proposition: universal relative reduced and operator reduced quotient}
    Let $f\colon P\to Q$ be a surjective quandle homomorphism.
    \begin{enumerate}
        \item Put $\redQuot{f}\coloneqq P/{\redCong{f}}$ and let $\redMap{f}\colon P\twoheadrightarrow\redQuot{f}$ be the quotient map.
        Then, the induced homomorphism $\redFac{f}\colon\redQuot{f}\twoheadrightarrow Q$ is reduced.
        
        Moreover, the factorization $f= \redFac{f}\circ \redMap{f}$ is universal among factorizations of $f$ whose second homomorphism is reduced.
        That is, if $f=h\circ u$ with $u\colon P\twoheadrightarrow R$ surjective and $h\colon R\twoheadrightarrow Q$ reduced, then $u$ factors uniquely through $\redQuot{f}$.

        \item Put $\opredQuot{f}\coloneqq P/\opredDef{f}$ and let $\opredMap{f}\colon P\twoheadrightarrow\opredQuot{f}$ be the quotient map.
        Then, the induced homomorphism $\opredFac{f}\colon\opredQuot{f}\twoheadrightarrow Q$ is operator-reduced.

        Moreover, the factorization $f= \opredFac{f}\circ \opredMap{f}$ is universal among factorizations of $f$ whose second homomorphism is operator-reduced.
        That is, if $f=h\circ u$ with $u\colon P\twoheadrightarrow R$ surjective and $h\colon R\twoheadrightarrow Q$ operator-reduced, then $u$ factors uniquely through $\opredQuot{f}$.
    \end{enumerate}
    \[\begin{tikzcd}[column sep=small]
        P \arrow[two heads]{rr}{f} \arrow[two heads]{rd}[swap]{\redMap{f}} & & Q\rlap{,} \\
        & \redQuot{f} \arrow[two heads]{ru}[swap]{\redFac{f}}
    \end{tikzcd}\qquad
    \begin{tikzcd}[column sep=small]
        P \arrow[two heads]{rr}{f} \arrow[two heads]{rd}[swap]{\opredMap{f}} & & Q\rlap{.} \\
        & \opredQuot{f} \arrow[two heads]{ru}[swap]{\opredFac{f}}
    \end{tikzcd}\]
\end{prop}

\begin{proof}
    Let $\pi=\redMap{f}$ and put $\bar f=\redFac{f}$. By \cref{lemma: relative symmetry groups under a surjective factorization}, every $\psi\in\relInn(\bar f)$ has the form $\pi_*(\varphi)$ for some $\varphi\in\relInn(f)$. Hence, $s_{\psi([x])}([x])=[s_{\varphi(x)}(x)]=[x]$ by the definition of $\redCong{f}$. Thus, $\redFac{f}$ is reduced.

    The same argument for $\pi=\opredMap{f}$ gives $[s_{[x]},s_{\psi([x])}] =\pi_*([s_x,s_{\varphi(x)}])=1$, since $[s_x,s_{\varphi(x)}]\in\opredDef{f}$. Hence, $\opredFac{f}$ is operator-reduced.

    By \cref{lemma: relative symmetry groups under a surjective factorization}, the defining relations of $\redCong{f}$ and the generators of $\opredDef{f}$ vanish in the corresponding quotient factors. The same lemma shows that they vanish in every reduced, respectively operator-reduced, quotient factorization of $f$. This proves the two universal properties. 
\end{proof}

\begin{dfn}
\label[dfn]{definition: reduced and operator-reduced factor}
    For a surjective quandle homomorphism $f\colon P \twoheadrightarrow Q$, the induced homomorphisms $\redFac{f}\colon\redQuot{f}\twoheadrightarrow Q$ and $\opredFac{f}\colon\opredQuot{f}\twoheadrightarrow Q$ are called the \emph{reduced factor} and the \emph{operator-reduced factor} of $f$, respectively.
\end{dfn}

\begin{lem}
\label[lem]{proposition: reduction preserves connected components}
    Let $f\colon P \twoheadrightarrow Q$ be a surjective quandle homomorphism.
    The quotient map $\redMap{f}\colon P\to\redQuot{f}$ induces a bijection $\pi_0(P)\to\pi_0(\redQuot{f})$.
\end{lem}

\begin{proof}
Since $\redMap{f}$ is surjective, $\pi_0(\redMap{f})$ is surjective.
Each defining pair of $\redCong{f}$ has the same image in $\pi_0(P)$. Hence, the canonical homomorphism $P\twoheadrightarrow\pi_0(P)$ factors through $\redMap{f}$ as
\[\begin{tikzcd}
    P \arrow[two heads]{r} \arrow[two heads]{d}[swap]{\redMap{f}} & \pi_0(P)\rlap{.} \\
    \redQuot{f} \arrow[dashed]{ru}[swap]{h}
\end{tikzcd}\]
Moreover, as $\pi_0(P)$ is a trivial quandle, $h$ factors through the canonical map $\redQuot{f} \to \pi_0(\redQuot{f})$. Writing $\overline{h}\colon \pi_0(\redQuot{f}) \to \pi_0(P)$ for the induced homomorphism, we have $\overline{h}\circ \pi_0(\redMap{f}) = \id$ by the universality, which shows that $\pi_0(\redMap{f})$ is also injective.
Thus, it is bijective.
\end{proof}


We conclude this subsection with further properties of the reduced factor. First, the reduced congruence has the following orbit description. Put $\redOpGrp_f(x)\coloneqq \langle s_{\varphi(x)}\mid\varphi\in\relInn(f)\rangle$.

\begin{prop}
\label[prop]{proposition: orbit description of relative reduction}
Let $f\colon P\twoheadrightarrow Q$ be an operator-reduced quandle homomorphism.
Then, $\redOpGrp_f(x)$ is abelian for every $x\in P$, and the fibers
of $\redMap{f}\colon P\twoheadrightarrow\redQuot{f}$ are the orbits
$\redOpGrp_f(x)\cdot x$.
\end{prop}

\begin{proof}
For $\varphi,\psi\in\relInn(f)$, operator-reducedness gives
\[ [s_{\varphi(x)},s_{\psi(x)}] = \varphi[s_x,s_{\varphi^{-1}\psi(x)}]\varphi^{-1} = 1. \]
Thus, $\redOpGrp_f(x)$ is abelian.
We note that $\theta \redOpGrp_f(x) \theta^{-1}= \redOpGrp_f(\theta(x))$ for $\theta\in \Inn(P)$, because $\theta s_{\varphi(x)} \theta^{-1} = s_{\theta\varphi(x)} = s_{\theta\varphi\theta^{-1}(\theta x)}$.

We define $x \sim_{\redOpGrp_f} y$ precisely when $y \in \redOpGrp_f(x)\cdot x$. We next want to show that the relations $\sim_{\redOpGrp_f}$ and $\redCong{f}$ coincide.

We observe that $\sim_{\redOpGrp_f}$ is a quandle congruence. Since $s_x\in \redOpGrp_f(x)$, we have $x\in \redOpGrp_f(x)\cdot x$, and hence $x \sim_{\redOpGrp_f} x$.
If $x\sim_{\redOpGrp_f} y$, then $y=b(x)$ for some $b \in \redOpGrp_f(x)$. As $\redOpGrp_f(x)$ is abelian, we have $\redOpGrp_f(y) = b\redOpGrp_f(x) b^{-1} = \redOpGrp_f(x) bb^{-1}=\redOpGrp_f(x)$. Hence, $x=b^{-1}(y)$ and $b^{-1}\in \redOpGrp_f(y)$ imply $x \in \redOpGrp_f(y)\cdot y$, which means $y \sim_{\redOpGrp_f} x$.
If $x \sim_{\redOpGrp_f} y$ and $y \sim_{\redOpGrp_f} z$, then a similar argument yields that $\redOpGrp_f(x)= \redOpGrp_f(y) = \redOpGrp_f(z)$ and $x\sim_{\redOpGrp_f}z$. Thus, $\sim_{\redOpGrp_f}$ is an equivalence relation.
Moreover, if $x \sim_{\redOpGrp_f} x'$ and $y \sim_{\redOpGrp_f} y'$, then we see $s_x=s_{x'}$. Therefore, as $\sim_{\redOpGrp_f}$ is preserved by $\Inn(P)$, we obtain $s_x^\varepsilon(y)\sim s_x^\varepsilon(y') =s_{x'}^\varepsilon(y')$ for $\varepsilon\in\{\pm 1\}$. Hence, $\sim_{\redOpGrp_f}$ is indeed a quandle congruence.

For each defining pair $x \redCong{f} s_{\varphi(x)}(x)$, we have $x \sim_{\redOpGrp_f} s_{\varphi(x)}(x)$ since $s_{\varphi(x)}(x) \in \redOpGrp_f(x)\cdot x$. Thus, $x\redCong{f}y$ implies $x\sim_{\redOpGrp_f} y$.
Conversely, if $y=s_{\varphi(x)}(x)$ for some $\varphi \in\relInn(f)$, then $x\redCong{f} s_{\varphi(x)}(x)=y$. Since the operators $s_{\varphi(x)}$ generate $\redOpGrp_f(x)$, we see that if $x\sim_{\redOpGrp_f} y$, then $x\redCong{f}y$.
Thus, the two congruences coincide.
\end{proof}






\subsection{Nilpotency ascent and commutative action criterion}
\label{subsection: quadratic relative actions and nilpotency ascent}

In this subsection, we study sufficient conditions for operator-reduced homomorphisms to be nilpotent.

The operator-reduced condition controls the action on abelian sections of the relative inner automorphism group. Using this, we first prove that an operator-reduced homomorphism with a finite relative inner automorphism group is nilpotent whenever its target is nilpotent.

For this, we need to make an observation on operator-reduced homomorphisms. 
Recall that the \emph{Fitting subgroup} $\Fit(G)$ of a group $G$ is the subgroup generated by all normal nilpotent subgroups of $G$. When $G$ is finite, $\Fit(G)$ is a finite nilpotent normal subgroup of $G$.

\begin{prop}
\label[prop]{proposition: fitting control for operator reduced homomorphisms}
Let $f\colon P\twoheadrightarrow Q$ be an operator-reduced quandle homomorphism, and let $M\subseteq \Inn(P)$ be a finite normal subgroup with $M\subseteq \relInn(f)$. Then, $[\Inn(P), M]\subseteq\Fit(M)$.
%Equivalently, $M/\Fit(M)\subseteq Z(G/\Fit(M))$. In particular, $M$ is solvable.
\end{prop}

\begin{proof}
Let $x\in P$, and put $B_x\coloneqq\redOpGrp_f(x)\cap M$, where $\redOpGrp_f(x)$ is the abelian group defined in \cref{proposition: orbit description of relative reduction}. Since $M\subseteq \relInn(f)$, we have $m\redOpGrp_f(x)m^{-1} =\redOpGrp_f(m(x))=\redOpGrp_f(x)$ for every $m\in M$. Thus, $B_x$ is an abelian normal subgroup of $M$, and hence $B_x\subseteq\Fit(M)$.

For $m\in M$, we see $[m,s_x]=s_{m(x)}s_x^{-1}\in B_x\subseteq \Fit(M)$. 
Since the symmetry maps $s_x$ generate $\Inn(P)$ and the Fitting group $\Fit(M)$ is normal in $\Inn(P)$, it follows that $[M, \Inn(P)]=[\Inn(P), M]\subseteq\Fit(M)$.
%Finally, $\Fit(M)$ is nilpotent and $M/\Fit(M)$ is abelian, so $M$ is solvable.
\end{proof}

The following elementary group-theoretic observation will be useful.

\begin{lem}
\label[lem]{lemma: quadratic action on abelian sections}
    Let $M\subseteq M'$ be normal subgroups of a group $G$ such that $A\coloneqq M'/M$ is abelian. Let $s\in G$ be an element satisfying $[[x,s],s]=1$ for all $x\in M'$. Then, letting $\alpha(s)\in\Aut(A)$ be induced by conjugation by $s$, we have $(\alpha(s)-\id_A)^2=0$.
\end{lem}

\begin{proof}
Let $x\in M'$, and write $\overline{x}=xM\in A$. Writing $A$ additively, the image of $[x,s]=xsx^{-1}s^{-1}$ in $A$ is $\overline{x}-\alpha(s)(\overline{x})=(\id_A-\alpha(s))(\overline{x})$.
Therefore $[[x,s],s]=1$ yields $(\id_A-\alpha(s))^2(\overline{x})=0$ for every $x\in M'$. Hence we have $(\alpha(s)-\id_A)^2=0$ in $\End(A)$.
\end{proof}


\begin{lem}[{\cite[Lemma 3.11]{book_isaacs_2008_finite_group_theory}}]
\label{fact:finite_solvable_minimal_normal_subgroup_is_elementary_abelian_p-group}
    Let $M$ be a solvable minimal normal subgroup of an arbitrary group $G$.
    Then $M$ is abelian, and if $M$ is finite, it is an elementary abelian $p$-group for some prime $p$.
\end{lem}

\begin{lem}% standard fact
\label{lemma:finite_nilpotent_generated_by_p-elements_is_p-group}
    A finite nilpotent group generated by $p$-elements is a $p$-group.
\end{lem}

\begin{proof}
    In a finite nilpotent group, every Sylow $p$-group is normal, and hence unique. Since $p$-elements are contained in the unique Sylow $p$-group, the generating condition implies that the finite nilpotent group itself is a $p$-group.
\end{proof}

\begin{lem}
\label{lemma:invariant_space_of_finite_vector_space_with_p-group_action_is_nontrivial}
    Let $H$ be a finite $p$-group, and let $V$ be a nonzero finite-dimensional vector space over $\mathbb{F}_{p}$.
    Suppose that $H$ acts linearly on $V$.
    Then the fixed-point subspace $V^H=\{v \in V \mid h\cdot v=v \text{ for all } h\in H \}$ is nonzero.
\end{lem}

\begin{proof}
    For each $v \in V$, the orbit $H\cdot v$ has cardinality dividing $\lvert H \rvert$. Since $H$ is a $p$-group, $\lvert H\cdot v\rvert$ is a power of $p$. In particular, every nontrivial orbit has cardinality divisible by $p$. The elements of $V^H$ are precisely those lying in orbits of cardinality $1$. Hence, by the orbit-counting formula,
    \[ \lvert V \rvert = \lvert V^H \rvert + (\text{a multiple of }p). \]
    Since $\lvert V\rvert$ is also divisible by $p$, we obtain $\lvert V^H \rvert \equiv 0 \mod p$. As $0\in V^H$, we have $\lvert V^H \rvert \geq 1$, and hence in fact $\lvert V^H \rvert \geq p$. Thus $V^H\neq 0$. 
\end{proof}

Recall that a \emph{$G$-chief series of $N$} for a normal subgroup $N\trianglelefteq G$ is a finite chain $1= N_{0}< N_{1}<\cdots< N_{m}=N$ of normal subgroups of $G$ such that $N_i/N_{i-1}$ is a nontrivial minimal normal subgroup of $G/N_{i-1}$ for every $1\leq i \leq m$.

\begin{thm}
\label{theorem:nilpotency ascent theorem for groups}
    Let $G$ be a group, $N\leq G$ be a finite normal subgroup, and $F\leq G$ be a finite nilpotent normal subgroup with $[G,N]\subseteq F\subseteq N$.
    Suppose that $G$ is generated by a set $S\subseteq G$ and that $[[x,s],s]=1$ for every $s\in S$ and $x\in N$.
    If $G/N$ is nilpotent, then $G$ is nilpotent.
    More precisely, if $G/N$ is nilpotent of class at most $r$ and $N$ has a $G$-chief series of length $m$, then $N$ is $G$-nilpotent of class at most $m$ and $G$ is nilpotent of class at most $m+r$.
\end{thm}

\begin{proof}
    \textbf{Step 1:}
    % N: finite solvable and G/F: nilpotent
    By assumption, $[N,N] \subseteq [G,N] \subseteq F \subseteq N$, so $N/F$ is abelian. Since $F$ is nilpotent, it is solvable, and hence $N$ is solvable. Moreover, $[G,N] \subseteq F$ implies $N/F \subseteq Z(G/F)$. Since $(G/F)/(N/F)\cong G/N$ is nilpotent, $G/F$ is nilpotent.

    \textbf{Step 2:}
    % Each V_i=N_i/N_{i-1} is a vec sp with linear action of $\alpha_i(G)$
    Consider the upper central series $1=Z_0(F)\leq Z_1(F)\leq\cdots\leq Z_c(F)=F$ of $F$. Since $F\trianglelefteq G$, each $Z_j(F)$ is normal in $G$. Thus
    \begin{equation}
        1=Z_0(F) \leq Z_1(F) \leq \cdots \leq Z_c(F)=F \leq N \label{equation:G-normal_series_of_N}
    \end{equation}
    is a $G$-normal series for $N$. Since $N$ is finite, this series can be refined to a $G$-chief series
    \[ 1=N_0 < N_1 < \cdots < N_m=N. \]
    For each \(1\le i\le m\), put $V_i=N_i/N_{i-1}$. Since $N$ is finite and solvable, each $V_i$ is finite and solvable. Moreover, $V_i$ is a minimal normal subgroup of $G/N_{i-1}$, so by \cref{fact:finite_solvable_minimal_normal_subgroup_is_elementary_abelian_p-group}, $V_i$ is an elementary abelian $p_i$-group for some prime $p_i$. Thus $V_i$ is a finite-dimensional vector space over $\mathbb{F}_{p_i}$, and $\Aut(V_i)=\GL_{\mathbb{F}_{p_i}}(V_i)$. 
    The conjugation action of $G$ on $V_i$ gives a homomorphism
    \[ \alpha_i\colon G\longrightarrow \Aut(V_i), \quad \alpha_i(g)(v)= gvg^{-1}. \]
    Hence $V_i$ naturally admits a linear action of $\alpha_i(G)$.

    \textbf{Step 3:}
    % \alpha_i(F)=1
    Fix $1\leq i \leq m$. Since the chosen $G$-chief series refines the series \eqref{equation:G-normal_series_of_N} we have either $N_i \leq F$ or $F\leq N_{i-1}$.
    If $N_i \leq F$, then there is some $j$ such that $Z_{j-1}(F) \leq N_{i-1} \leq N_i \leq Z_j(F)$. Hence $[N_i, F] \leq [Z_j(F), F] \leq Z_{j-1}(F) \leq N_{i-1}$.
    If $F\leq N_{i-1}$, then $[N_i, F] \leq F \leq N_{i-1}$ since $F \trianglelefteq G$.
    Thus, in either case, we obtain $[N_i,F] \leq N_{i-1}$, which means $[F,V_i]=1$, or, equivalently, $\alpha_i(F)=1$.

    \textbf{Step 4:}
    % \alpha_i(G): finite nilpotent
    Since $\alpha_i(F)=1$, the homomorphism $\alpha_i\colon G \to \Aut(V_i)$ factors through $G/F$. In particular, there is a surjective homomorphism $G/F \twoheadrightarrow \alpha_i(G)$.
    By Step 1, $G/F$ is nilpotent, so $\alpha_i(G)$ is also nilpotent. Moreover, since $V_i$ is finite, so $\Aut(V_i)$, hence $\alpha_i(G)$, is finite. Thus $\alpha_i(G)$ is a finite nilpotent group.

    \textbf{Step 5:}
    % \alpha_i(G): finite p-group
    By assumption, $[[x,s],s]=1$ for $x\in N$ and $s\in S$. By passing to the factor $V_i$ and writing additively, we obtain $(\alpha_i(s)-1)^2=0$ in 
    $\End(V_i)$
    %$\Aut(V_i)=\GL_{\mathbb{F}_{p_i}}(V_i)$ 
    by \cref{lemma: quadratic action on abelian sections}.
    Putting $u=\alpha_i(s)-1$, we have $u^2=0$. Since $V_i$ is a vector space over $\mathbb{F}_{p_i}$, its endomorphism ring has characteristic $p_i$. Hence we have
    \[ (\alpha_i(s))^{p_i} = (u+1)^{p_i} = u^{p_i} + 1 =1, \]
    because $u^2=0$ implies $u^{p_i}=0$. Thus every $\alpha_i(s)$ is a $p_i$-element. Since $\alpha_i(G)$ is a finite nilpotent group generated by the $p_i$-elements $\alpha_i(s)$, it follows that $\alpha_i(G)$ is a $p_i$-group (\cref{lemma:finite_nilpotent_generated_by_p-elements_is_p-group}).

    \textbf{Step 6:}
    % \alpha_i(G)=1 
    Thus the finite $p_i$-group $\alpha_i(G)$ acts linearly on the nonzero finite-dimensional $\mathbb{F}_{p_i}$-vector space $V_i$. By \cref{lemma:invariant_space_of_finite_vector_space_with_p-group_action_is_nontrivial}, $V_i^{\alpha_i(G)}\neq 0$. Since the action is induced by conjugation, we have
    \[ V_i^{\alpha_i(G)}= \{v \in V_i \mid \alpha_i(g)(v)=gvg^{-1}=v \text{ for all } g\in G \} = Z(G/N_{i-1},V_i)\neq 0, \]
    which is a normal subgroup of $G/N_{i-1}$. By the minimality of $V_i$, we obtain $V_i^{\alpha_i(G)}=V_i$. This means $\alpha_i(G)=1$. 

    \textbf{Step 7:}
    % N: G-nilpotent and G: nilpotent
    Since $\alpha_i(G)=1$, we have $[G,N_i] \subseteq N_{i-1}$ for every $1\leq i \leq m$. The chosen $G$-chief series of $N$
    \[ 1=N_0 \leq N_1 \leq \cdots \leq N_m=N \]
    is in fact a $G$-central series, and hence $N$ is $G$-nilpotent. Finally, since $G/N$ is nilpotent, it follows that $G$ is nilpotent.
\end{proof}

\begin{thm}[Nilpotency ascent theorem]
\label[thm]{theorem: finite operator reduced ascent}
    Let $f\colon P\twoheadrightarrow Q$ be an operator-reduced quandle homomorphism such that $\relInn(f)$ is a finite group. If $Q$ is nilpotent, then both $f$ and $P$ are nilpotent. 
    More precisely, if $Q$ is $r$-nilpotent and $\relInn(f)$ has an $\Inn(P)$-chief series of length $m$, then $f$ is $m$-nilpotent and $P$ is $(m+r)$-nilpotent.
\end{thm}

\begin{proof}
This follows immediately from \cref{proposition: commutator interpretation of reducedness,proposition: fitting control for operator reduced homomorphisms} and \cref{theorem:nilpotency ascent theorem for groups}, by taking $G=\Inn(P)$, $N=\relInn(f)$, and $F=\Fit(N)$.
\end{proof}



\begin{cor}
\label[cor]{corollary: consequences of finite ascent}
    Suppose that $f$ is operator-reduced and that $\relInn(f)$ is finite. Then $P$ is nilpotent if and only if $Q$ is nilpotent. Under these equivalent conditions, $f$ is nilpotent. 
\end{cor}

\begin{proof}
If $Q$ is nilpotent, the assertion follows from \cref{theorem: finite operator reduced ascent}. Conversely, if $P$ is nilpotent, then $Q$ is nilpotent because $\Inn(Q)$ is a quotient of $\Inn(P)$. The last assertion follows from \cref{theorem: finite operator reduced ascent}.
\end{proof}

\begin{rem}
Note that if $f$ is reduced, or if $f$ is connected and every fiber of $f$ is trivial, then $f$ is operator-reduced. Thus, the same equivalence holds in these cases.
\end{rem}


The absolute theorem of Darn\'e is recovered without any additional work: if $P$ is finitely generated and reduced, then $P$ is nilpotent by \cite[Proposition~7.13]{darne_2026_nilpotent_quandles}, and hence every surjective homomorphism out of $P$ is nilpotent by \cref{corollary: relative nilpotency implies ordinary nilpotency}. The following finite example shows why no direct relative analogue can hold without a condition on the base.

\begin{eg}%[Necessity of the nilpotent-base hypothesis]
\label[eg]{example: finite connected reduced homomorphism not nilpotent}
Let $T_{n}$ be the subquandle of $\Conj(S_n)$ consisting of
the transpositions. The quotient
$S_4\twoheadrightarrow S_4/V_4\cong S_3$ induces a surjective
homomorphism
$f\colon T_{4}\twoheadrightarrow T_{3}$.
Since the transpositions generate $S_n$ and $Z(S_n)=1$ for $n=3,4$,
we have
$\Inn(T_{4})\cong S_4$ and
$\Inn(T_{3})\cong S_3$, and
$\relInn(f)=V_4$. The fibers are
$\{(12),(34)\}$, $\{(13),(24)\}$, and $\{(14),(23)\}$.
The group $V_4$ acts transitively on each fiber, so $f$ is connected.
Elements in the same fiber are equal or disjoint transpositions, and
therefore commute. Hence every fiber is trivial and $f$ is reduced.

On the other hand, $[S_4,V_4]=V_4$, so
$\relLCS_{i}(S_4,V_4)=V_4$ for every $i\geq0$. Thus $f$ is not
nilpotent.
\end{eg}

This example is finite, connected, has trivial fibers, and has an abelian
relative inner automorphism group. In particular, finite generation of
the domain, finite generation of the kernel congruence, or finite
normal generation of the relative inner automorphism group cannot
replace the nilpotency hypothesis on the target in
\cref{theorem: finite operator reduced ascent}. 
%The obstruction is visible in \cref{proposition: finite relative nilpotency via chief factors}: $V_4$ is a minimal nontrivial normal subgroup of $S_4$, hence a $S_4$-chief factor, but it is not central. Equivalently, the induced action of $S_4$ on $V_4$ is nontrivial.



The second result is the following commutative action criterion. Unlike the preceding ascent theorem, the following quantitative criterion does not require the relative inner automorphism group to be finite.

\begin{lem}
\label[lem]{lemma: abelian images of inner automorphism groups}
Let $P$ be a quandle satisfying either of the following conditions:
\begin{enumerate}[label={(\alph*)}]
    \item\label{item:Inn(P)_is_finitely_generated} $\Inn(P)$ is generated by finitely many symmetry maps. 
    \item\label{item:P_has_finite_connected_component} $P$ has finitely many connected components.
\end{enumerate}
Let $\beta\colon\Inn(P)\to H$ be a group homomorphism with abelian image.
Then, the image $\beta(\Inn(P))$ is generated by finitely many elements of the form $\beta(s_{x_1}),\ldots,\beta(s_{x_r})$.
\end{lem}

\begin{proof}
For the first case \ref{item:Inn(P)_is_finitely_generated}, if $\Inn(P)$ is generated by $s_{x_1},\dots,s_{x_r}$ for $x_i\in P$, then $\beta(\Inn(P))$ is generated by $\beta(s_{x_1}),\ldots,\beta(s_{x_r})$.

For the second case \ref{item:P_has_finite_connected_component}, suppose that $P$ has $r<\infty$ connected components and choose $x_1,\ldots,x_r\in P$ from each connected component.
For every $x\in P$, there are some $i\in \{1,\dots,r\}$ and $\varphi\in\Inn(P)$ such that $x=\varphi(x_i)$. Since $s_x=\varphi s_{x_i}\varphi^{-1}$ and $\beta(\Inn(P))$ is abelian, we have
\[ \beta(s_x) = \beta(\varphi)\beta(s_{x_i})\beta(\varphi)^{-1} = \beta(s_{x_i}). \]
Since $\Inn(P)$ is generated by all symmetry maps, its image under $\beta$ is generated by $\beta(s_{x_1}),\ldots,\beta(s_{x_r})$.
\end{proof}

\begin{thm}[Commutative action criterion]
\label[thm]{theorem: commutative action criterion}
    Let $f\colon P\twoheadrightarrow Q$ be an operator-reduced quandle homomorphism. 
    Assume that there is a finite sequence of $\Inn(P)$-normal subgroups $1=N_0\subseteq N_1\subseteq\cdots\subseteq N_d=\relInn(f)$ such that each $A_j=N_j/N_{j-1}$ is abelian and the image of the conjugation action $\alpha_j\colon \Inn(P)\to \Aut(A_j)$ is abelian and generated by finitely many elements of the form $\alpha_j(s_{x_1}),\dots,\alpha_j(s_{x_r})$ for $x_i\in P$. Then, $f$ is nilpotent of class at most $d(r+1)$.
\end{thm}

\begin{proof}
Let $G=\Inn(P)$ and $N=\relInn(f)$.
For a fixed $j$, put $\alpha_i=\alpha_j(s_{x_i})$ for $1\leq i\leq r$. From the assumption, $\alpha_j(G)$ is an abelian group generated by $\alpha_1,\ldots,\alpha_r$.

Let $R_j$ be the subring of $\End(A_j)$ generated by $\alpha_1^{\pm1},\ldots,\alpha_r^{\pm1}$, which is in fact a commutative ring by the assumption. Put $u_i=\alpha_i-\id$ and $I_j=(u_1,\ldots,u_r)\subseteq R_j$. By \cref{proposition: commutator interpretation of reducedness,lemma: quadratic action on abelian sections}, we have $u_i^2=0$. Since the $u_i$ commute, every product of $r+1$ of them contains a repeated factor. Hence we obtain $I_j^{r+1}=0$.

Since $\alpha_i^{-1}=\id-u_i$, $\alpha_i^{\pm1} \equiv\id \pmod{I_j}$. Since $\alpha_j(G)$ is generated by $\alpha_1,\ldots,\alpha_r$, every image $\alpha_j(g)$ is congruent to $\id$ modulo $I_j$, and hence the action difference $\id - \alpha_j(g)$ belongs to $I_j$.
Because taking the commutator of an element in $A_j$ with an element in $G$ is equivalent to multiplying it by an element in $I_j$ of the form $\id - \alpha_j(g)$, the $(r+1)$-fold commutator of $A_j$ with $G$ is trivial. In the language of subgroups, this implies
\[  [N_j, \underbrace{G, \dots, G}_{r+1}] \subseteq N_{j-1}.\]
Applying this relation successively along the series $1 = N_0 \subseteq N_1 \subseteq \cdots \subseteq N_d = N$ of length $d$, we deduce that $N$ is annihilated by $d(r+1)$ commutators with $G$. This immediately yields $N\subseteq Z_{d(r+1)}(G)$, and hence $\relInn(f)$ is nilpotent of class at most $d(r+1)$ by \cref{theorem: characterizations of nilpotent groups}. 
\end{proof}



\begin{cor}
\label[cor]{corollary: easy consequences of comm action criterion}
    Let $f\colon P\twoheadrightarrow Q$ be an operator-reduced quandle homomorphism. 
    Suppose that $P$ satisfies either of the two conditions in \cref{lemma: abelian images of inner automorphism groups}, and let $r$ denote the number of symmetry maps needed there to generate the corresponding abelian image.
    If $\relInn(f)$ is nilpotent of class at most $d$ and $\Inn(Q)$ is abelian, then $f$ is $d(r+1)$-nilpotent.
    %In particular, if both $\relInn(f)$ and $\Inn(Q)$ are abelian, then $f$ is $(r+1)$-nilpotent.
\end{cor}

\begin{proof}
    Put $G=\Inn(P)$ and $N=\relInn(f)$. Consider the upper central series of \(N\):
    \[ 1=Z_0(N)\leq Z_1(N)\leq\cdots\leq Z_d(N)=N. \]
    Since each \(Z_j(N)\) is characteristic in \(N\) and \(N\trianglelefteq G\), this is a \(G\)-normal series. Set $ A_j=Z_j(N)/Z_{j-1}(N)$ for $1\leq j\leq d$. Each \(A_j\) is abelian. Moreover, by the definition of the upper central series, $[N,Z_j(N)]\subseteq Z_{j-1}(N)$, which means that \(N\) acts trivially on \(A_j\) by conjugation. It follows that the action of \(G\) on \(A_j\) factors through $G/N\cong \Inn(Q)$.
    Since \(\Inn(Q)\) is abelian, the image of \(G\) in \(\Aut(A_j)\) is abelian. Thus all the hypotheses of the commutative action criterion are satisfied. Consequently, \cref{theorem: commutative action criterion} shows that $f$ is $d(r+1)$-nilpotent.
\end{proof}





If $P$ is generated by $r$ elements as a quandle, then \cref{lemma: quandle generators generate inner group} shows that $\Inn(P)$ is generated by the corresponding $r$ symmetry maps.


Finally, we study the cyclic case.

\begin{prop}
\label[prop]{proposition: operator reduced homomorphisms with cyclic relative inner group}
Let $f\colon P\twoheadrightarrow Q$ be a surjective quandle homomorphism such that $\relInn(f)$ is cyclic. Then, $f$ is operator-reduced if and
only if it is $2$-nilpotent.
Under these equivalent conditions, $f$ is $1$-nilpotent if $\relInn(f)$ is infinite or has square-free order.
\end{prop}

\begin{proof}
The implication from $2$-nilpotency to operator-reducedness follows from \cref{proposition: commutator interpretation of reducedness}. Suppose that $f$ is operator-reduced.
Let $G=\Inn(P)$.
Put $N=\relInn(f)$, write it additively, and for $x\in P$, let $\alpha_x\in\Aut(N)$ be induced by conjugation by $s_x$. Put $u_x\coloneqq\alpha_x-\id_N$. Since $[[a,s_x],s_x]=1$ for every $a\in N$, \cref{lemma: quadratic action on abelian sections} gives $u_x^2=0$.

If $N$ is infinite cyclic, then $\End(N)\cong\ZZ$, so $u_x=0$ for every $x\in P$. Since the symmetry maps generate $G$, we have $[G,N]=1$, and hence $f$ is $1$-nilpotent.

Suppose that $N\cong\ZZ/n\ZZ$ for some $n\geq1$, and write $\textstyle n=\prod_{p\mid n}p^{e_p}$. Under $\End(N)\cong\ZZ/n\ZZ$, each $u_x$ is multiplication by an integer $a_x$. The equality $u_x^2=0$ means that $n\mid a_x^2$, and hence $p^{\lceil e_p/2\rceil}\mid a_x$ for every prime $p\mid n$. Put $\textstyle k\coloneqq\prod_{p\mid n}p^{\lceil e_p/2\rceil}$. Then, $k\mid a_x$ and $n\mid k^2$, so $u_x(N)\subseteq kN$ and $u_x(kN)=0$ for every $x\in P$. Thus, every symmetry map acts trivially on both $N/kN$ and $kN$. Since the symmetry maps generate $G$, it follows that $[G,N]\subseteq kN$ and $[G,kN]=1$. Therefore,
\[
    \relInnLCS{2}{f} = [G,[G,N]] \subseteq [G,kN] = 1.
\]
Hence, $f$ is $2$-nilpotent. If $n$ is square-free, then $k=n$ and $kN=0$, so $[G,N]=1$. Thus, $f$ is $1$-nilpotent in this case as well.
\end{proof}
